\documentclass[12pt]{article}
\usepackage{amsmath, amssymb, geometry}
\usepackage{hyperref}
\geometry{a4paper, margin=1in}

\title{Homework 6.2: Homogeneous Linear Equations with Constant Coefficients \\ \large Comprehensive Problem Sets \& Solutions}
\author{}
\date{}

\begin{document}

\maketitle
\tableofcontents
\newpage

\section{Concept Checklist \& Question Mapping}
This section outlines all the core concepts, problem types, and algebraic tricks covered in the source material. Each concept is mapped to specific problems in the generated problem set to ensure 100\% conceptual coverage.

\begin{itemize}
    \item \textbf{Concept 1: Distinct Real Roots} (Basic 2nd, 3rd, and 4th order equations) \\
    \textit{Questions: 1, 2, 3}
    \item \textbf{Concept 2: Repeated Real Roots} (Handling multiplicity and the "Missing $x$" trap) \\
    \textit{Questions: 4, 5, 6}
    \item \textbf{Concept 3: Complex Conjugate Roots} (Euler's formula, identifying $\alpha$ and $\beta$) \\
    \textit{Questions: 7, 8, 9}
    \item \textbf{Concept 4: Repeated Complex Roots} (Multiplying by $x$ for repeated imaginary roots) \\
    \textit{Questions: 10, 11}
    \item \textbf{Concept 5: Operator Form} (Extracting roots directly from differential operators) \\
    \textit{Questions: 12, 13}
    \item \textbf{Concept 6: Initial Value Problems (IVPs)} (Applying product/chain rules for constants) \\
    \textit{Questions: 14, 15, 16, 17}
    \item \textbf{Concept 7: Systems of Differential Equations} (Algebraic elimination method) \\
    \textit{Questions: 18, 19, 20}
    \item \textbf{Concept 8: Cauchy-Euler Equations} (Variable coefficients, $y=x^r$ substitution) \\
    \textit{Questions: 21, 22, 23}
    \item \textbf{Concept 9: Boundary Value Problems (BVPs)} (Conditions at two different points) \\
    \textit{Questions: 24, 25}
    \item \textbf{Concept 10: Higher-Order \& Algebraic Tricks} (Factoring by grouping, quadratic form, synthetic division) \\
    \textit{Questions: 26, 27, 28}
    \item \textbf{Concept 11: Real-World Modeling} (Vibrating beam, Momentum Gradient Descent) \\
    \textit{Questions: 29, 30}
\end{itemize}

\newpage

\section{Problem Sets}
\textbf{Instructions:} Find the general solution for the following differential equations unless initial/boundary conditions are provided. 

\subsection*{Distinct Real Roots}
\textbf{1.} $y'' - 5y' + 6y = 0$ \\
\textbf{2.} $y''' - 2y'' - y' + 2y = 0$ \\
\textbf{3.} $y^{(4)} - 5y'' + 4y = 0$

\subsection*{Repeated Real Roots}
\textbf{4.} $y'' - 8y' + 16y = 0$ \\
\textbf{5.} $y''' + 3y'' + 3y' + y = 0$ \\
\textbf{6.} $y''' - y'' = 0$

\subsection*{Complex Conjugate Roots}
\textbf{7.} $y'' + 9y = 0$ \\
\textbf{8.} $y'' - 4y' + 13y = 0$ \\
\textbf{9.} $y''' - y'' + 4y' - 4y = 0$

\subsection*{Repeated Complex Roots}
\textbf{10.} $y^{(4)} + 8y'' + 16y = 0$ \\
\textbf{11.} $(D^2 - 2D + 5)^2[y] = 0$

\subsection*{Operator Form}
\textbf{12.} $(D-2)^3(D+1)[y] = 0$ \\
\textbf{13.} $D^2(D^2+9)^2[y] = 0$

\subsection*{Initial Value Problems (IVPs)}
\textbf{14.} $y'' - 3y' - 4y = 0; \quad y(0)=1, \quad y'(0)=6$ \\
\textbf{15.} $y'' + 4y' + 4y = 0; \quad y(0)=2, \quad y'(0)=-1$ \\
\textbf{16.} $y'' + 2y' + 2y = 0; \quad y(0)=0, \quad y'(0)=1$ \\
\textbf{17.} $y''' - y' = 0; \quad y(0)=1, \quad y'(0)=0, \quad y''(0)=1$

\subsection*{Systems of Differential Equations}
\textbf{18.} Solve the system: $\frac{dx}{dt} = 3x - y, \quad \frac{dy}{dt} = 9x - 3y$ \\
\textbf{19.} Solve the system: $x' = x + y, \quad y' = 4x + y$ \\
\textbf{20.} Solve the system: $x'' + y' = 0, \quad x' + y'' = 0$

\subsection*{Cauchy-Euler Equations ($x > 0$)}
\textbf{21.} $x^2y'' - 2xy' + 2y = 0$ \\
\textbf{22.} $x^2y'' + 5xy' + 4y = 0$ \\
\textbf{23.} $x^2y'' - xy' + 5y = 0$

\subsection*{Boundary Value Problems (BVPs)}
\textbf{24.} $y'' + 16y = 0; \quad y(0)=0, \quad y(\pi/2)=0$ \\
\textbf{25.} $y'' + 16y = 0; \quad y(0)=0, \quad y(\pi/8)=1$

\subsection*{Higher-Order \& Algebraic Tricks}
\textbf{26.} $y''' - 5y'' - y' + 5y = 0$ (Hint: Factor by grouping) \\
\textbf{27.} $y^{(4)} - 8y'' - 9y = 0$ (Hint: Quadratic form) \\
\textbf{28.} $y''' - 4y'' + 5y' - 2y = 0$ (Hint: Rational Root Theorem / Synthetic Division)

\subsection*{Real-World Modeling}
\textbf{29.} \textbf{Vibrating Beam:} The transverse displacement of a uniform beam satisfies $y^{(4)} - 16y = 0$. Find the general solution. \\
\textbf{30.} \textbf{Momentum Gradient Descent:} A machine learning parameter $\theta(t)$ undergoing continuous-time momentum gradient descent on a quadratic loss surface is modeled by $\ddot{\theta} + 6\dot{\theta} + 9\theta = 0$. Find the general solution and determine if the system is underdamped, overdamped, or critically damped.

\newpage

\section{Detailed Solutions}

\subsection*{Distinct Real Roots}
\textbf{Solution 1:} 
Auxiliary equation: $r^2 - 5r + 6 = 0$. 
Factoring yields $(r-2)(r-3) = 0$. 
Roots are $r_1 = 2, r_2 = 3$. 
General solution: $y(x) = c_1e^{2x} + c_2e^{3x}$.

\textbf{Solution 2:} 
Auxiliary equation: $r^3 - 2r^2 - r + 2 = 0$. 
Factor by grouping: $r^2(r-2) - 1(r-2) = 0 \implies (r^2-1)(r-2) = 0 \implies (r-1)(r+1)(r-2) = 0$. 
Roots are $r = 1, -1, 2$. 
General solution: $y(x) = c_1e^x + c_2e^{-x} + c_3e^{2x}$.

\textbf{Solution 3:} 
Auxiliary equation: $r^4 - 5r^2 + 4 = 0$. 
This is quadratic in form. Let $u = r^2$, then $u^2 - 5u + 4 = 0 \implies (u-1)(u-4) = 0$. 
Substituting back: $(r^2-1)(r^2-4) = 0 \implies (r-1)(r+1)(r-2)(r+2) = 0$. 
Roots are $r = \pm 1, \pm 2$. 
General solution: $y(x) = c_1e^x + c_2e^{-x} + c_3e^{2x} + c_4e^{-2x}$.

\subsection*{Repeated Real Roots}
\textbf{Solution 4:} 
Auxiliary equation: $r^2 - 8r + 16 = 0$. 
Factoring yields $(r-4)^2 = 0$. 
Root is $r = 4$ with multiplicity 2. 
General solution: $y(x) = c_1e^{4x} + c_2xe^{4x}$.

\textbf{Solution 5:} 
Auxiliary equation: $r^3 + 3r^2 + 3r + 1 = 0$. 
Recognize the perfect cube binomial expansion: $(r+1)^3 = 0$. 
Root is $r = -1$ with multiplicity 3. 
General solution: $y(x) = c_1e^{-x} + c_2xe^{-x} + c_3x^2e^{-x}$.

\textbf{Solution 6:} 
Auxiliary equation: $r^3 - r^2 = 0$. 
Factor out $r^2$: $r^2(r-1) = 0$. 
Roots are $r = 0$ (multiplicity 2) and $r = 1$. 
General solution: $y(x) = c_1e^{0x} + c_2xe^{0x} + c_3e^x = c_1 + c_2x + c_3e^x$.

\subsection*{Complex Conjugate Roots}
\textbf{Solution 7:} 
Auxiliary equation: $r^2 + 9 = 0 \implies r^2 = -9 \implies r = \pm 3i$. 
Here $\alpha = 0, \beta = 3$. 
General solution: $y(x) = c_1\cos(3x) + c_2\sin(3x)$.

\textbf{Solution 8:} 
Auxiliary equation: $r^2 - 4r + 13 = 0$. 
Using the quadratic formula: $r = \frac{4 \pm \sqrt{16 - 52}}{2} = \frac{4 \pm \sqrt{-36}}{2} = \frac{4 \pm 6i}{2} = 2 \pm 3i$. 
Here $\alpha = 2, \beta = 3$. 
General solution: $y(x) = e^{2x}(c_1\cos(3x) + c_2\sin(3x))$.

\textbf{Solution 9:} 
Auxiliary equation: $r^3 - r^2 + 4r - 4 = 0$. 
Factor by grouping: $r^2(r-1) + 4(r-1) = 0 \implies (r^2+4)(r-1) = 0$. 
Roots are $r = 1$ and $r = \pm 2i$. 
General solution: $y(x) = c_1e^x + c_2\cos(2x) + c_3\sin(2x)$.

\subsection*{Repeated Complex Roots}
\textbf{Solution 10:} 
Auxiliary equation: $r^4 + 8r^2 + 16 = 0$. 
Quadratic in form: $(r^2+4)^2 = 0$. 
Roots are $r = \pm 2i$ with multiplicity 2. 
General solution: $y(x) = c_1\cos(2x) + c_2\sin(2x) + c_3x\cos(2x) + c_4x\sin(2x)$.

\textbf{Solution 11:} 
Auxiliary equation: $(r^2 - 2r + 5)^2 = 0$. 
Solve the inner quadratic: $r = \frac{2 \pm \sqrt{4-20}}{2} = 1 \pm 2i$. 
Because of the squared factor, the roots $1 \pm 2i$ have multiplicity 2. 
General solution: $y(x) = e^x(c_1\cos(2x) + c_2\sin(2x)) + xe^x(c_3\cos(2x) + c_4\sin(2x))$.

\subsection*{Operator Form}
\textbf{Solution 12:} 
The operator form translates directly to the auxiliary equation: $(r-2)^3(r+1) = 0$. 
Roots are $r = 2$ (multiplicity 3) and $r = -1$. 
General solution: $y(x) = c_1e^{2x} + c_2xe^{2x} + c_3x^2e^{2x} + c_4e^{-x}$.

\textbf{Solution 13:} 
Auxiliary equation: $r^2(r^2+9)^2 = 0$. 
Roots are $r = 0$ (multiplicity 2) and $r = \pm 3i$ (multiplicity 2). 
General solution: $y(x) = c_1 + c_2x + c_3\cos(3x) + c_4\sin(3x) + c_5x\cos(3x) + c_6x\sin(3x)$.

\subsection*{Initial Value Problems (IVPs)}
\textbf{Solution 14:} 
Auxiliary equation: $r^2 - 3r - 4 = 0 \implies (r-4)(r+1) = 0 \implies r = 4, -1$. 
General solution: $y(x) = c_1e^{4x} + c_2e^{-x}$. 
Derivative: $y'(x) = 4c_1e^{4x} - c_2e^{-x}$. 
Apply conditions: 
$y(0) = c_1 + c_2 = 1$ 
$y'(0) = 4c_1 - c_2 = 6$. 
Add the equations: $5c_1 = 7 \implies c_1 = 7/5$. 
Substitute back: $7/5 + c_2 = 1 \implies c_2 = -2/5$. 
Particular solution: $y(x) = \frac{7}{5}e^{4x} - \frac{2}{5}e^{-x}$.

\textbf{Solution 15:} 
Auxiliary equation: $r^2 + 4r + 4 = 0 \implies (r+2)^2 = 0 \implies r = -2$ (mult 2). 
General solution: $y(x) = c_1e^{-2x} + c_2xe^{-2x}$. 
Apply $y(0) = 2 \implies c_1 = 2$. 
Derivative (use product rule): $y'(x) = -2c_1e^{-2x} + c_2(e^{-2x} - 2xe^{-2x})$. 
Apply $y'(0) = -1 \implies -2(2) + c_2(1) = -1 \implies -4 + c_2 = -1 \implies c_2 = 3$. 
Particular solution: $y(x) = 2e^{-2x} + 3xe^{-2x}$.

\textbf{Solution 16:} 
Auxiliary equation: $r^2 + 2r + 2 = 0 \implies r = \frac{-2 \pm \sqrt{4-8}}{2} = -1 \pm i$. 
General solution: $y(x) = e^{-x}(c_1\cos x + c_2\sin x)$. 
Apply $y(0) = 0 \implies 1 \cdot (c_1(1) + c_2(0)) = 0 \implies c_1 = 0$. 
So, $y(x) = c_2e^{-x}\sin x$. 
Derivative: $y'(x) = c_2(-e^{-x}\sin x + e^{-x}\cos x)$. 
Apply $y'(0) = 1 \implies c_2(0 + 1) = 1 \implies c_2 = 1$. 
Particular solution: $y(x) = e^{-x}\sin x$.

\textbf{Solution 17:} 
Auxiliary equation: $r^3 - r = 0 \implies r(r^2-1) = 0 \implies r = 0, 1, -1$. 
General solution: $y(x) = c_1 + c_2e^x + c_3e^{-x}$. 
Derivatives: $y'(x) = c_2e^x - c_3e^{-x}$, $y''(x) = c_2e^x + c_3e^{-x}$. 
Apply conditions: 
$y(0) = c_1 + c_2 + c_3 = 1$ 
$y'(0) = c_2 - c_3 = 0 \implies c_2 = c_3$ 
$y''(0) = c_2 + c_3 = 1 \implies 2c_2 = 1 \implies c_2 = 1/2, c_3 = 1/2$. 
Substitute into first eq: $c_1 + 1/2 + 1/2 = 1 \implies c_1 = 0$. 
Particular solution: $y(x) = \frac{1}{2}e^x + \frac{1}{2}e^{-x}$ (which is $\cosh x$).

\subsection*{Systems of Differential Equations}
\textbf{Solution 18:} 
From eq 1, isolate $y$: $y = 3x - x'$. 
Differentiate: $y' = 3x' - x''$. 
Substitute $y$ and $y'$ into eq 2: 
$3x' - x'' = 9x - 3(3x - x') \implies 3x' - x'' = 9x - 9x + 3x' \implies -x'' = 0 \implies x'' = 0$. 
Auxiliary equation for $x$: $r^2 = 0 \implies r = 0$ (mult 2). 
Solution for $x$: $x(t) = c_1 + c_2t$. 
Find $y$: $y(t) = 3(c_1 + c_2t) - (c_2) = (3c_1 - c_2) + 3c_2t$.

\textbf{Solution 19:} 
From eq 1, isolate $y$: $y = x' - x$. 
Differentiate: $y' = x'' - x'$. 
Substitute into eq 2: 
$x'' - x' = 4x + (x' - x) \implies x'' - 2x' - 3x = 0$. 
Auxiliary equation: $r^2 - 2r - 3 = 0 \implies (r-3)(r+1) = 0 \implies r = 3, -1$. 
Solution for $x$: $x(t) = c_1e^{3t} + c_2e^{-t}$. 
Find $y$: $y(t) = x' - x = (3c_1e^{3t} - c_2e^{-t}) - (c_1e^{3t} + c_2e^{-t}) = 2c_1e^{3t} - 2c_2e^{-t}$.

\textbf{Solution 20:} 
From eq 1, $y' = -x''$. Differentiate to get $y'' = -x'''$. 
Substitute into eq 2: $x' - x''' = 0 \implies x''' - x' = 0$. 
Auxiliary equation: $r^3 - r = 0 \implies r(r^2-1) = 0 \implies r = 0, 1, -1$. 
Solution for $x$: $x(t) = c_1 + c_2e^t + c_3e^{-t}$. 
Find $y$: We know $y' = -x'' = -(c_2e^t + c_3e^{-t})$. 
Integrate to find $y$: $y(t) = -c_2e^t + c_3e^{-t} + c_4$. 
(Check with eq 2: $x' + y'' = (c_2e^t - c_3e^{-t}) + (-c_2e^t + c_3e^{-t}) = 0$. It holds). 
General solution: $x(t) = c_1 + c_2e^t + c_3e^{-t}$, $y(t) = c_4 - c_2e^t + c_3e^{-t}$.

\subsection*{Cauchy-Euler Equations}
\textbf{Solution 21:} 
Substitute $y = x^r$. The characteristic equation is $r(r-1) - 2r + 2 = 0$. 
Expand: $r^2 - 3r + 2 = 0 \implies (r-1)(r-2) = 0 \implies r = 1, 2$. 
General solution: $y(x) = c_1x + c_2x^2$.

\textbf{Solution 22:} 
Substitute $y = x^r$. Characteristic equation: $r(r-1) + 5r + 4 = 0$. 
Expand: $r^2 + 4r + 4 = 0 \implies (r+2)^2 = 0 \implies r = -2$ (mult 2). 
General solution: $y(x) = c_1x^{-2} + c_2x^{-2}\ln x$.

\textbf{Solution 23:} 
Substitute $y = x^r$. Characteristic equation: $r(r-1) - r + 5 = 0$. 
Expand: $r^2 - 2r + 5 = 0$. 
Roots: $r = \frac{2 \pm \sqrt{4-20}}{2} = 1 \pm 2i$. 
General solution: $y(x) = x(c_1\cos(2\ln x) + c_2\sin(2\ln x))$.

\subsection*{Boundary Value Problems (BVPs)}
\textbf{Solution 24:} 
Auxiliary equation: $r^2 + 16 = 0 \implies r = \pm 4i$. 
General solution: $y(x) = c_1\cos(4x) + c_2\sin(4x)$. 
Apply $y(0) = 0 \implies c_1(1) + c_2(0) = 0 \implies c_1 = 0$. 
So $y(x) = c_2\sin(4x)$. 
Apply $y(\pi/2) = 0 \implies c_2\sin(2\pi) = 0 \implies c_2(0) = 0$. 
This is true for any $c_2$. 
Conclusion: There are infinitely many solutions, $y(x) = c_2\sin(4x)$.

\textbf{Solution 25:} 
From the previous problem, general solution is $y(x) = c_1\cos(4x) + c_2\sin(4x)$. 
Apply $y(0) = 0 \implies c_1 = 0$. 
So $y(x) = c_2\sin(4x)$. 
Apply $y(\pi/8) = 1 \implies c_2\sin(4(\pi/8)) = 1 \implies c_2\sin(\pi/2) = 1 \implies c_2(1) = 1 \implies c_2 = 1$. 
Conclusion: There is a unique solution, $y(x) = \sin(4x)$.

\subsection*{Higher-Order \& Algebraic Tricks}
\textbf{Solution 26:} 
Auxiliary equation: $r^3 - 5r^2 - r + 5 = 0$. 
Factor by grouping: $r^2(r-5) - 1(r-5) = 0 \implies (r^2-1)(r-5) = 0 \implies (r-1)(r+1)(r-5) = 0$. 
Roots: $r = 1, -1, 5$. 
General solution: $y(x) = c_1e^x + c_2e^{-x} + c_3e^{5x}$.

\textbf{Solution 27:} 
Auxiliary equation: $r^4 - 8r^2 - 9 = 0$. 
Quadratic form: $(r^2-9)(r^2+1) = 0 \implies (r-3)(r+3)(r^2+1) = 0$. 
Roots: $r = \pm 3, \pm i$. 
General solution: $y(x) = c_1e^{3x} + c_2e^{-3x} + c_3\cos x + c_4\sin x$.

\textbf{Solution 28:} 
Auxiliary equation: $r^3 - 4r^2 + 5r - 2 = 0$. 
Using the Rational Root Theorem, test factors of $-2$: $\pm 1, \pm 2$. 
Test $r=1$: $1^3 - 4(1)^2 + 5(1) - 2 = 1 - 4 + 5 - 2 = 0$. So $r=1$ is a root. 
Divide polynomial by $(r-1)$ using synthetic division or polynomial long division to get $r^2 - 3r + 2$. 
Factor the quotient: $r^2 - 3r + 2 = (r-1)(r-2)$. 
Full factorization: $(r-1)^2(r-2) = 0$. 
Roots: $r = 1$ (mult 2), $r = 2$. 
General solution: $y(x) = c_1e^x + c_2xe^x + c_3e^{2x}$.

\subsection*{Real-World Modeling}
\textbf{Solution 29:} 
Auxiliary equation: $r^4 - 16 = 0$. 
Factor as difference of squares: $(r^2-4)(r^2+4) = 0 \implies (r-2)(r+2)(r^2+4) = 0$. 
Roots: $r = \pm 2, \pm 2i$. 
General solution: $y(x) = c_1e^{2x} + c_2e^{-2x} + c_3\cos(2x) + c_4\sin(2x)$.

\textbf{Solution 30:} 
Auxiliary equation: $r^2 + 6r + 9 = 0$. 
Factor: $(r+3)^2 = 0$. 
Roots: $r = -3$ (multiplicity 2). 
General solution: $\theta(t) = c_1e^{-3t} + c_2te^{-3t}$. 
Because the roots are real and repeated, the system is \textbf{critically damped}.

\newpage

\section{Concept Checklist \& Question Mapping (Reference)}
\begin{itemize}
    \item \textbf{Concept 1: Distinct Real Roots} $\rightarrow$ Questions: 1, 2, 3
    \item \textbf{Concept 2: Repeated Real Roots} $\rightarrow$ Questions: 4, 5, 6
    \item \textbf{Concept 3: Complex Conjugate Roots} $\rightarrow$ Questions: 7, 8, 9
    \item \textbf{Concept 4: Repeated Complex Roots} $\rightarrow$ Questions: 10, 11
    \item \textbf{Concept 5: Operator Form} $\rightarrow$ Questions: 12, 13
    \item \textbf{Concept 6: Initial Value Problems (IVPs)} $\rightarrow$ Questions: 14, 15, 16, 17
    \item \textbf{Concept 7: Systems of Differential Equations} $\rightarrow$ Questions: 18, 19, 20
    \item \textbf{Concept 8: Cauchy-Euler Equations} $\rightarrow$ Questions: 21, 22, 23
    \item \textbf{Concept 9: Boundary Value Problems (BVPs)} $\rightarrow$ Questions: 24, 25
    \item \textbf{Concept 10: Higher-Order \& Algebraic Tricks} $\rightarrow$ Questions: 26, 27, 28
    \item \textbf{Concept 11: Real-World Modeling} $\rightarrow$ Questions: 29, 30
\end{itemize}

\end{document}